% TODO proof-reading
\subsection{導入}

\newcommand{\JADURL}{\url{http://varaneckas.com/jad/}}

Java（または一般的に\ac{JVM}バイトコード）用のよく知られた逆コンパイラがいくつかあります
\footnote{例えば、JAD: \JADURL}。

その理由は、\ac{JVM}バイトコードの逆コンパイルが、
より低レベルなx86コードよりもやや簡単であることです：

\begin{itemize}
\item データ型に関する情報がはるかに多い。
\item \ac{JVM}のメモリモデルははるかに厳密で明確に定義されている。
\item Javaコンパイラは最適化を行わない（\ac{JVM} \ac{JIT}が実行時に行う）ため、
      クラスファイル内のバイトコードは通常かなり読みやすい。
      
\end{itemize}

\ac{JVM}の知識はいつ役立つでしょうか？

\newcommand{\URLListOfJVMLangs}{\url{http://en.wikipedia.org/wiki/List_of_JVM_languages}}

\begin{itemize}
\item 逆コンパイラの結果を再コンパイルする必要なしに、クラスファイルの手軽で簡単なパッチタスク。
\item 難読化されたコードの分析。
\item 独自の難読化ツールの構築。
\item \ac{JVM}をターゲットとするコンパイラのコードジェネレータ（バックエンド）の構築（Scala、Clojureなど
      \footnote{完全なリスト：\URLListOfJVMLangs}）。
      
\end{itemize}

いくつかの簡単なコード片から始めましょう。
特に記載がない限り、JDK 1.7を使用しています。

クラスファイルを逆コンパイルするために使用するコマンドです：\\
\GTT{javap -c -verbose}。

すべての例を準備する際に使用した本：\JavaBook。